% Lösungen Einheit 4 von 5 – Trapez, Drachenviereck, Raute (zusammenbau v0.1)
\einheitenkopf{Einheit 4 von 5 · Trapez, Drachenviereck, Raute}

%% TODO Lösung der Erklärzeile 22g) fehlt in der Bank
\erg{22}{a) $A = \frac{a + c}{2} \cdot h$ \quad b) $A = (9 + 5) : 2 \cdot 4 = 28$ cm² \quad c) $A = (14 + 8) : 2 \cdot 6 = 66$ m² \quad d) $A = (11 + 7) : 2 \cdot 5 = 45$ cm² \quad e) $A = (25 + 15) : 2 \cdot 9 = 180$ m² \quad f) $A = (16 + 6) : 2 \cdot 7 = 77$ dm² \quad g) \ldots \quad h) $A = (7{,}5 + 4{,}5) : 2 \cdot 3 = 18$ cm² \quad i) $h = 2 \cdot 42 : (9 + 5) = 6$ cm \quad j) $A = 7 \cdot 6 : 2 = 21$ cm² \quad k) $u = 11 + 5 + 2 \cdot 5 = 26$ cm (die Höhe braucht man nicht) \quad l) $h = 2 \cdot 2\,960 : (90 + 70) = 5\,920 : 160 = 37$ cm; $37$ cm $\ge 35$ cm – ja, die Töpfe passen.}
\erg{23}{a) Überstand $(14 - 8) : 2 = 3$ cm; Schenkel $\sqrt{3^2 + 4^2} = 5$ cm; $u = 14 + 8 + 2 \cdot 5 = 32$ cm}
\erg{24}{a) $f = 2 \cdot 35 : 7 = 10$ cm}
\erg{25}{a) Leo nimmt nur die Grundseite $a$, ohne Mittelwert der parallelen Seiten. Richtig: $A = (15 + 9) : 2 \cdot 7 = 84$ cm² \quad b) Legt man zwei gleiche Trapeze gedreht aneinander, entsteht ein Parallelogramm mit der Grundseite $a + c$ und der Höhe $h$. Ein Trapez ist die Hälfte davon. \quad c) $A = (32 + 6) : 2 \cdot 7 = 133$ m²; $133 \cdot 100 = 13\,300$ m³}
